\customSection[sec:design]{Проєктування та реалізація вебзастосунку}{0em}

\customSubSection[sec:style]{Обґрунтування архітектурного стилю}{1em}

Архітектурне рішення приймається відповідно до вимог розділу~I і описується за практиками, узагальненими в ISO/IEC/IEEE~42010:2022~\citecustom{17} та~\citecustom{18}: кожне рішення пов'язується з чинниками якості, а компроміси фіксуються явно. Розглянуто три кандидати: клієнтський застосунок, у якому вся логіка виконується в браузері; класичний клієнт-серверний застосунок із серверним розрахунковим ядром і базою даних; гібрид із локальним редактором і сервером розрахунків.

Профіль системи визначає вибір. Дані проєкту невеликі (кімнати, десятки елементів, кілька проводів), розрахунки виконуються за мілісекунди (підрозділ~\ref{sec:verification}), користувачі не потребують спільної роботи над проєктом, а вимоги NFR-04 і NFR-05 передбачають відсутність сервера та передавання даних. Основні ризики пов'язані не з масштабом, а з наочністю тривимірної моделі, відгуком інтерфейсу під час перетягування об'єктів і коректністю розрахунків.

\begin{itemize}[noitemsep, topsep=0.3em, leftmargin=*]
    \item \textbf{Клієнт-серверний варіант} додав би автентифікацію, базу даних, мережеві запити на кожне перетягування та супровід інфраструктури без користі для задачі; крім того, дані проєкту покидали б пристрій користувача, що суперечить NFR-05.
    \item \textbf{Гібридний варіант} доцільний, якщо розрахункове ядро потребує бібліотек, відсутніх у браузері. Алгоритми цього застосунку (геометрія, ортогональні маршрути, розкрій) виконуються на JavaScript за частки секунди, тому виокремлення сервера не виправдане.
    \item \textbf{Клієнтський варіант} дає миттєвий відгук, нульову вартість розгортання (статичні файли) та повну приватність даних; його недоліки~-- залежність від можливостей браузера та відсутність спільного доступу~-- для поставленої задачі несуттєві.
\end{itemize}

Обрано \textbf{клієнтський застосунок без збірки}: статичні файли HTML, CSS та ES-модулі~\citecustom{22}, графічна бібліотека Three.js~\citecustom{20}, підключена через карту імпортів~\citecustom{23}. Основні компроміси: (а)~відсутність сервера спрощує розгортання, але проєкт зберігається лише у файлі користувача; (б)~відсутність збірника спрощує супровід, але залишає завантаження Three.js із мережевого CDN, тож офлайн-робота потребує локальної копії бібліотеки; (в)~уся логіка в браузері обмежена його можливостями (зокрема, перезапис відкритого файлу на місці працює лише в браузерах на основі Chromium~\citecustom{24}).

\customSubSection[sec:stack]{Технологічний стек і склад програмного продукту}{1em}

Перелік обраних технологій наведено в \tabref{tab:stack}.

\begin{spacing}{1.15}
\small\noindent
\begin{xltabular}{\textwidth}{| >{\raggedright\arraybackslash}p{4.4cm} | >{\raggedright\arraybackslash}X |}
    \caption{Технологічний стек застосунку} \label{tab:stack} \\
    \hline
    \textbf{Складова} & \textbf{Технологія} \\ \hline
    \endfirsthead

    \hline
    \textbf{Складова} & \textbf{Технологія} \\ \hline
    \endhead

    Мова та модулі & JavaScript, ES-модулі браузера~\citecustom{22} \\ \hline
    Тривимірна графіка & Three.js 0.160~\citecustom{20}, \texttt{OrbitControls}; рендеринг у WebGL~\citecustom{21} \\ \hline
    Підключення бібліотек & карта імпортів (\texttt{importmap}) із CDN unpkg~\citecustom{23} \\ \hline
    Розмітка та стилі & HTML5, CSS3 (змінні CSS, темна тема) \\ \hline
    Збереження даних & формат JSON~\citecustom{25}, File System Access API та завантаження файлу~\citecustom{24} \\ \hline
    Розгортання & статичний HTTP-сервер (наприклад, \texttt{python -m http.server})~\citecustom{26} \\ \hline
    Збірка, база даних, сервер застосунку & не використовуються \\ \hline
\end{xltabular}
\end{spacing}

Мова JavaScript обрана як основна мова браузерного середовища, що виконується без плагінів, а Three.js надає високорівневий доступ до WebGL~\citecustom{21}: сцену, камеру, матеріали, трасування променів для вибору об'єктів. Модулі ES дозволяють розділити код без збірника, але вимагають завантаження сторінки за протоколом HTTP, а не з локальної файлової системи, тому застосунок запускається статичним сервером. Карта імпортів зіставляє ім'я \texttt{three} з адресою бібліотеки в мережі.

Склад програмного продукту наведено в \tabref{tab:modules}. Алгоритм розкрою ГКЛ винесено в окремий модуль, який не залежить ні від DOM, ні від Three.js: він отримує описи поверхонь як числа й повертає план, тому його можна запускати й перевіряти окремо (NFR-03).

\begin{spacing}{1.15}
\small\noindent
\begin{xltabular}{\textwidth}{| >{\raggedright\arraybackslash}p{3.3cm} | >{\raggedright\arraybackslash}X | c |}
    \caption{Склад програмного продукту} \label{tab:modules} \\
    \hline
    \textbf{Файл} & \textbf{Призначення} & \textbf{Рядків} \\ \hline
    \endfirsthead

    \hline
    \textbf{Файл} & \textbf{Призначення} & \textbf{Рядків} \\ \hline
    \endhead

    \texttt{index.html} & розмітка, чотири вкладки, карта імпортів & 607 \\ \hline
    \texttt{app.js} & оркестрація: стан, сцена Three.js, форми, взаємодія, збереження, розрахунки, трасування & 3694 \\ \hline
    \texttt{js/drywall.js} & чисте ядро розкрою ГКЛ (розбиття, пул відрізків, оптимізація) & 503 \\ \hline
    \texttt{js/utils.js} & допоміжні функції (\texttt{round2}, \texttt{hexToNum}, \texttt{colorHex}, \texttt{clampv}) & 15 \\ \hline
    \texttt{style.css} & тема інтерфейсу & 650 \\ \hline
\end{xltabular}
\end{spacing}

\customSubSection[sec:c4]{Архітектурні подання за моделлю C4}{1em}

Архітектуру описано за нотацією C4~\citecustom{19}, що передбачає чотири ієрархічні рівні абстракції~-- контекст, контейнери, компоненти та код.

\textbf{Рівень 1 (контекст).} Система взаємодіє з одним типом користувача (проєктувальник, виконроб, власник житла) та двома зовнішніми сутностями: мережею доставки бібліотек (CDN) і локальною файловою системою, у якій зберігаються файли проєктів. Зовнішніх інформаційних систем, від яких залежить розрахунок, немає.

\textbf{Рівень 2 (контейнери)} показано на рис.~\ref{fig:c4}. Єдиним виконуваним контейнером є вебклієнт у браузері; статичний HTTP-сервер лише доставляє файли.

\settikz[fig:c4]{\textwidth}{assets/c4-containers}{Діаграма C4 рівня контейнерів}

\textbf{Рівень 3 (компоненти \texttt{app.js}).} Усередині оркестратора виділяються функціональні групи: \emph{стан} (об'єкт \texttt{state}, знімок початкового стану, історія); \emph{геометрія} (контекст кімнати, базис стіни, перетворення «зсув від кута~$\to$ світові координати»); \emph{побудова сцени} (\texttt{buildScene}, меші стін, дверей, елементів, світильників, предметів, проводів, накладка розкладки); \emph{форми й списки} (кімнати, двері, електрообладнання, проводи, світильники, предмети); \emph{розрахунки} (площі, ГКЛ, профілі, плінтуси, довжина проводів); \emph{взаємодія} (вибір об'єктів променем, перетягування, створення проводів, гарячі клавіші); \emph{збереження} (серіалізація, імпорт, нормалізація, скасування). Єдине тривимірне полотно переміщується між вкладками «Конструктор» і «Проводка», а на інших вкладках розміщується в прихованому контейнері.

\customSubSection[sec:data]{Модель даних і збереження проєкту}{1em}

Стан застосунку зберігається в одному об'єкті \texttt{state}. Його основні сутності та зв'язки показано на рис.~\ref{fig:model}. Кожна кімната має прямокутний контур у плані, центр у світових координатах, висоту, рівень підлоги й товщину кожної з чотирьох стін. Двері, елементи, світильники та предмети посилаються на кімнату ідентифікатором. Провід посилається на два вузли (елемент або світильник) і містить необов'язковий список контрольних точок, кожна з яких є вільною точкою $(x,y,z)$ або посиланням на вузол.

\settikz[fig:model]{0.98\textwidth}{assets/er-diagram}{Концептуальна модель даних проєкту}

Проєкт серіалізується у файл JSON~\citecustom{25} версії~2 із верхніми ключами: \texttt{app}, \texttt{version}, \texttt{rooms}, \texttt{activeRoom}, \texttt{room} (параметри відображення та розкладки), \texttt{doors}, \texttt{elements}, \texttt{wires}, \texttt{lights}, \texttt{wallObjects}. Зображення стін, дверей і картин зберігаються як data-URL усередині JSON, тому файл може значно збільшуватися.

Під час імпорту дані не приймаються як є, а нормалізуються: числові поля перевіряються на скінченність і додатність, перелічувані значення (стіна, тип, сторона відліку, режим розкладки) звіряються з дозволеними, невідомі кімнати замінюються першою, записи з хибними параметрами відкидаються, а проводи залишаються лише ті, обидва кінці яких існують. Передбачено сумісність із файлами версії~1, у яких описувалася єдина кімната. Така нормалізація забезпечує цілісність даних (NFR-07), але не є повною: наприклад, назви кімнат не перевіряються на наявність розмітки (підрозділ~\ref{sec:verification}).

Історія скасування зберігає серіалізовані знімки стану. Знімок записується перед змінами через форми та перед кожною взаємодією у 3D (перетягуванням, створенням проводу, додаванням чи видаленням контрольної точки); глибина історії обмежена 60 знімками, а скасування (\texttt{Ctrl+Z}) повторно застосовує знімок через той самий механізм імпорту без повторного кадрування камери.

\customSubSection[sec:geometry]{Геометрична модель приміщення}{1em}

Система координат правобічна, вісь $y$ спрямована вгору, вісь $x$ відповідає ширині кімнати, вісь $z$~-- довжині. Кожна кімната має центр $(c_x,c_z)$, ширину $w$, довжину $l$, висоту $h$ та рівень підлоги $F$. Чотири стіни названо за сторонами світу; для кожної визначено напрямок відліку від лівого кута (як для спостерігача всередині кімнати) і напрямок зовнішньої нормалі (\tabref{tab:walls}).

\begin{spacing}{1.15}
\small\noindent
\begin{xltabular}{\textwidth}{| l | c | >{\raggedright\arraybackslash}X | c |}
    \caption{Стіни кімнати та перетворення «зсув $a$ від лівого кута $\to$ світові координати»} \label{tab:walls} \\
    \hline
    \textbf{Стіна} & \textbf{Довжина} & \textbf{Точка на стіні $(x,z)$ на відстані $a$} & \textbf{Нормаль назовні} \\ \hline
    \endfirsthead

    \hline
    \textbf{Стіна} & \textbf{Довжина} & \textbf{Точка на стіні $(x,z)$ на відстані $a$} & \textbf{Нормаль назовні} \\ \hline
    \endhead

    north & $w$ & $(c_x-w/2+a,\; c_z+l/2)$ & $(0,+1)$ \\ \hline
    south & $w$ & $(c_x+w/2-a,\; c_z-l/2)$ & $(0,-1)$ \\ \hline
    east & $l$ & $(c_x-w/2,\; c_z-l/2+a)$ & $(-1,0)$ \\ \hline
    west & $l$ & $(c_x+w/2,\; c_z+l/2-a)$ & $(+1,0)$ \\ \hline
\end{xltabular}
\end{spacing}

Користувач вказує зсув від лівого або правого кута. Для елементів, світильників і предметів відстань від лівого кута дорівнює $a=o$ для лівого відліку та $a=L-o$ для правого, де $o$~-- зсув, $L$~-- довжина стіни. Для дверей зсув відліковується до краю прорізу, тому ліва межа прорізу дорівнює $o$ або $L-o-b$, де $b$~-- ширина прорізу. Товщина стіни відкладається назовні від внутрішнього контуру, тому внутрішні розміри кімнати від неї не залежать; якщо дві кімнати мають спільну стіну товщиною $t$, відстань між їхніми центрами дорівнює $w_1/2+t+w_2/2$ (у контрольній сцені $5/2+0{,}1+4/2=4{,}6$~м), а якщо кожна має власну стіну~-- $w_1/2+t_1+t_2+w_2/2$. Рівень підлоги кожної кімнати задається окремо, що дозволяє моделювати перепади висот.

Двері й арки подаються контуром на стіні, а за ненульових відкосів~-- трьома додатковими поверхнями (дві щоки та перемичка), що йдуть назовні на задану глибину. Двері мають полотно, яке обертається навколо петлі на заданий кут у бік приміщення або назовні, лиштву та поріг. Арка~-- прохід без полотна із заокругленим верхом.

\customSubSection[sec:wiring]{Електрообладнання та трасування проводки}{1em}

\customSubSubSection[sec:wiring-elements]{Типи елементів}{0.5em}

Електрообладнання розміщується на стінах, світильники~-- на стелі або стіні. Типи елементів та їхні розміри наведено в \tabref{tab:elements}; розмір щитка користувач може змінити.

\begin{spacing}{1.15}
\small\noindent
\begin{xltabular}{\textwidth}{| >{\raggedright\arraybackslash}p{5.2cm} | c | >{\raggedright\arraybackslash}X |}
    \caption{Типи електрообладнання} \label{tab:elements} \\
    \hline
    \textbf{Тип} & \textbf{Розмір, м} & \textbf{Особливість зображення} \\ \hline
    \endfirsthead

    \hline
    \textbf{Тип} & \textbf{Розмір, м} & \textbf{Особливість зображення} \\ \hline
    \endhead

    Розетка & $0{,}09\times0{,}09$ & один модуль \\ \hline
    Розетка подвійна & $0{,}17\times0{,}09$ & рамка з двома модулями \\ \hline
    Розетка потрійна & $0{,}25\times0{,}09$ & рамка з трьома модулями \\ \hline
    Вимикач & $0{,}09\times0{,}09$ & одна клавіша \\ \hline
    Вимикач подвійний / потрійний & $0{,}09\times0{,}09$ & дві / три клавіші \\ \hline
    Комбінований (розетка та вимикач) & $0{,}09\times0{,}17$ & розетка знизу, вимикач зверху \\ \hline
    Вивід проводки & $0{,}05\times0{,}05$ & пластина з петлею кабелю \\ \hline
    Електрощиток & $0{,}32\times0{,}50$ & розмір змінюється \\ \hline
    Розподільча коробка & $0{,}10\times0{,}10$ & квадратна коробка \\ \hline
\end{xltabular}
\end{spacing}

\customSubSubSection[sec:wiring-route]{Алгоритм ортогонального трасування}{0.5em}

Провід з'єднує два вузли~$P$ і~$Q$ (елементи чи світильники). Маршрут будується як ламана з прямими кутами; його схему показано на рис.~\ref{fig:route}. Вузол елемента розташований на стіні на висоті $F+e$ і виступає всередину кімнати на 0,03~м, де $e$~-- висота від підлоги; настінний світильник відсунуто від стіни на 0,1~м, а стельовий розташований на висоті $F+h-0{,}06$~м.

\settikz[fig:route]{0.72\textwidth}{assets/route-flow}{Схема алгоритму ортогонального трасування проводу}

Якщо обидва вузли лежать на одній стіні однієї кімнати й користувач не вимагає прокладання через стелю, маршрут має форму літери «Г»: горизонтальний відрізок на висоті вузла $P$, потім вертикальний до вузла $Q$. В іншому разі провід піднімається до рівня прокладання
\begin{equation}\label{eq:cy}
    c_y = \min\bigl(F_P+h_P,\; F_Q+h_Q\bigr) - 0{,}02,
\end{equation}
де $F_P, h_P$ і $F_Q, h_Q$~-- рівень підлоги та висота кімнат вузлів. Мінімум забезпечує, що горизонтальний перехід між кімнатами проходить під стелями обох. Далі провід рухається в горизонтальній площині за схемою Мангеттена~-- спершу вздовж осі $x$, потім вздовж осі $z$ (перегин у точці $(Q_x,c_y,P_z)$ додається лише тоді, коли обидві різниці ненульові) і опускається до вузла $Q$. Довжина маршруту
\begin{equation}\label{eq:wirelen}
    L = \underbrace{|c_y-P_y|}_{\text{підйом}} + |Q_x-P_x| + |Q_z-P_z| + \underbrace{|c_y-Q_y|}_{\text{спуск}},
\end{equation}
оскільки всі відрізки паралельні осям координат; для маршруту «Г» довжина дорівнює сумі горизонтального та вертикального відрізків. Контрольні точки розбивають провід на ділянки, і кожна ділянка трасується тим самим алгоритмом; точка може бути прив'язана до вузла, і тоді провід проходить через нього. Паралельні проводи зсуваються перпендикулярно напрямку на $0{,}05$~м на кожну «смугу», щоб не накладатися. Сумарна довжина із запасом на підключення $s$~\% дорівнює
\begin{equation}\label{eq:wiretotal}
    L_\Sigma = \Bigl(1+\frac{s}{100}\Bigr)\sum_i L_i .
\end{equation}

Для інтерактивного створення проводу користувач у режимі «Проводка» тягне курсор від одного вузла до іншого; під час руху показується тимчасова лінія. Якщо кнопку відпущено над іншим вузлом, провід створюється (дублікати відхиляються); якщо над існуючим проводом або його контрольною точкою~-- вузол вмикається в цей кабель як проміжна точка. Контрольну точку можна додати перетягуванням спеціальної кнопки на провід або кліком у режимі «Проводка», пересувати по стінах і стелі та видалити подвійним кліком.

\customSubSection[sec:calc]{Розрахунки матеріалів}{1em}

Розрахунки виконуються на вкладці «Розрахунки» одразу після зміни вхідних даних. Усі формули зведено нижче; параметри за замовчуванням наведено в дужках.

\textbf{Площі стін.} Для кімнати з розмірами $w\times l\times h$ повна площа стін і чиста площа після віднімання прорізів
\begin{equation}\label{eq:area}
    A_{\mathrm{gross}} = 2(w+l)\,h, \qquad A_{\mathrm{net}} = A_{\mathrm{gross}} - \sum_{d} b_d\, g_d ,
\end{equation}
де $b_d$ і $g_d$~-- ширина та висота $d$-го прорізу (дверей чи арки), що належить кімнаті. Результати округлюються до сотих метра квадратного.

\textbf{Відкоси та листи ГКЛ.} Площа відкосів прорізу з глибиною $r_d>0$ дорівнює $2 r_d g_d + r_d b_d$ (дві щоки та перемичка). Кількість листів розміру $s_w\times s_h$ (за замовчуванням $1{,}2\times2{,}5$~м) за відсотка відходів $p$ (10~\%)
\begin{equation}\label{eq:sheets}
    N = \left\lceil \frac{A\,(1+p/100)}{s_w\, s_h} \right\rceil , \qquad A = A_{\mathrm{net}} + A_{\mathrm{rev}} \;[+\,A_{\mathrm{ceil}}],
\end{equation}
де $A_{\mathrm{ceil}}=\sum w\,l$~-- площа стелі, яка додається за ввімкненого параметра.

\textbf{Профілі каркаса стелі.} Для кожної кімнати за кроку несучого профілю $\delta$ (0,6~м) число рядів $n=\lfloor l/\delta\rfloor+1$; довжина направляючого профілю UD дорівнює периметру $2(w+l)$, несучого CD~-- $n\,w$; кількість прямих підвісів $n(\lfloor w/0{,}6\rfloor+1)$, крабів~-- $n(\lfloor w/\delta\rfloor+1)$. Кількість планок профілю довжиною $\ell$ (3~м) дорівнює $\lceil L/\ell\rceil$.

\textbf{Плінтуси.} За сумарного периметра $P=\sum 2(w+l)$ та сумарної ширини прорізів $B=\sum b_d$ довжина підлогового плінтуса $\max(0,P-B)$, стельового~-- $P$; кількість планок довжиною $\ell$ (2,5~м) із запасом $p$ (5~\%) дорівнює $\lceil L\,(1+p/100)/\ell\rceil$. Двері з порогом додатково відмічаються окремим рядком у результатах.

\textbf{Довжина проводки} визначається за \eqref{eq:wirelen} і \eqref{eq:wiretotal}; результати виводяться окремо для кожного проводу з позначкою типу маршруту («по стіні», «по стелі», «через стелю», «власний маршрут») і сумарно (запас 15~\%).

Приклад. Для контрольної сцени (підрозділ~\ref{sec:verification}) чиста площа стін дорівнює $44{,}82+41{,}31=86{,}13$~м$^2$, площа відкосів $1{,}02+0{,}51+0{,}51=2{,}04$~м$^2$, тож $A=88{,}17$~м$^2$ і за \eqref{eq:sheets} $N=\lceil 88{,}17\cdot1{,}1/3{,}0\rceil=\lceil 32{,}33\rceil=33$ листи.

\customSubSection[sec:drywall]{Алгоритм розкрою гіпсокартону}{1em}

\customSubSubSection[sec:drywall-input]{Вхідні поверхні}{0.5em}

Планувальник отримує набір поверхонь, кожна з яких задана шириною~$W$, висотою~$H$ і списком прямокутних прорізів у координатах поверхні. Поверхнями є стіни кожної кімнати, відкоси дверей і арок (дві щоки та перемичка для кожного) та стелі. Прямокутні прорізи подаються габаритами двері чи арки; заокруглення верху арки в розкладці не враховується, тобто проріз розглядається як прямокутник.

\customSubSubSection[sec:drywall-main]{Основний режим}{0.5em}

Схему алгоритму наведено на рис.~\ref{fig:dwflow}.

\settikz[fig:dwflow]{0.9\textwidth}{assets/drywall-flow}{Схема алгоритму розкрою гіпсокартону}

\textbf{Розбиття на плитки.} Для кожної поверхні вибирається орієнтація листа: розмір плитки $s_w\times s_h$ (режим \texttt{h}) або $s_h\times s_w$ (режим \texttt{v}); у режимі «Авто» береться варіант із меншою кількістю плиток $\lceil W/c_w\rceil\lceil H/c_h\rceil$. Поверхня заповнюється рядами; за ввімкненого шахового порядку перша плитка кожного непарного ряду має половинну ширину, що зсуває вертикальні шви. Крайові плитки обрізаються до розміру поверхні.

\textbf{Врахування прорізів.} Для кожної плитки обчислюється площа перетину з прорізами; чиста площа плитки дорівнює площі плитки мінус площа перетину. Плитки, що повністю потрапили в проріз, пропускаються. Для частково перекритих плиток запам'ятовуються розміри вирізок, які пізніше потрапляють у пул відрізків.

\textbf{Пул відрізків.} Залишки зберігаються як прямокутники з позначкою листа-джерела. Відрізки, у яких одна зі сторін менша за 0,15~м або площа менша за 0,06~м$^2$, вважаються відходом і до пулу не потрапляють. Пошук виконується за найкращим збігом: серед відрізків, у які плитка вміщається (допускається поворот на 90°), береться той, що залишає найменшу площу. Після вирізання плитки відрізок ділиться гільйотинним розрізом на два прямокутники; напрямок розрізу обирається так, щоб більший із залишків залишився цілим. Перший розріз проводиться на повну висоту відрізка (залишки $r_w\times H$ і $p_w\times t_h$), якщо
\begin{equation}\label{eq:split}
    \max(r_w H,\; p_w t_h) \;\ge\; \max(r_w p_h,\; W t_h),
\end{equation}
і на повну ширину (залишки $r_w\times p_h$ і $W\times t_h$) в іншому разі; тут $W\times H$~-- розмір відрізка, $p_w\times p_h$~-- вирізана плитка, $r_w=W-p_w$ і $t_h=H-p_h$.

\textbf{Проходи.} У проході~1 цілі плити купуються як нові листи, а вирізки з їхніх отворів додаються до пулу. У проході~2 крайові плитки впорядковуються за спаданням площі (за аналогією з евристиками «перший підходящий за спаданням»~\citecustom{5}) і кожна або береться з пулу, або вимагає нового листа з додаванням залишку до пулу. У необов'язковому проході~3 (прапорець «Використати залишки на стінах») найменш використані куплені листи по черзі виключаються: лист разом з усіма його відрізками вилучається з пулу, і якщо його чисту площу можна скласти зі шматків пулу, лист скасовується, а плитка отримує статус «зі шматків». Складання виконує рекурсивний пошук (\texttt{coverRect}): одним відрізком, а якщо не вдається~-- гільйотинним розрізом за розмірами наявних відрізків; кількість операцій обмежено 60\,000 на спробу. Прохід повторюється, доки є покращення.

\customSubSubSection[sec:drywall-econ]{Режим «Економія»}{0.5em}

У режимі «Економія» прорізи вирізаються заздалегідь: кожна поверхня розкладається на прямокутні області чистої площі, які нарізаються на плитки, а плитки впорядковуються за спаданням площі й розміщуються через загальний пул, що зберігає навіть дрібні відрізки (від 4~см і 25~см$^2$). Випробовуються три стратегії нарізання (найкраща для кожної області, усі горизонтальні, усі вертикальні), і вибирається та, що купує найменше нових листів, а за рівності~-- з меншим залишком. Шви не зміщуються: це ціна мінімізації відходів.

\customSubSubSection[sec:drywall-compl]{Складність і результати}{0.5em}

Нехай $P$~-- кількість плиток, $K$~-- розмір пулу. Пошук відрізка для кожної плитки виконується за $O(K)$, тож проходи~1 і~2 працюють за $O(P\,K)$ із сортуванням $O(P\log P)$. Прохід~3 у найгіршому разі повторює пошук складання з обмеженням 60\,000 операцій на спробу, але на практиці вимірюваний час значно менший (підрозділ~\ref{sec:verification}). Алгоритм є евристикою: він не гарантує глобального мінімуму кількості листів, а лише наближається до нього; для оцінки якості використовується нижня межа \eqref{eq:lowerbound}. На виході планувальник повертає плитки з номером, статусом (ціла, ціла з отвором, розріз, із залишку, зі шматків), джерелом залишку, кількістю нових листів, повторних використань, зібраних зі шматків і список залишків.

\customSubSection[sec:ui]{Інтерфейс користувача}{1em}

Інтерфейс складається з чотирьох вкладок. Вкладка «Файл» містить створення нового проєкту, імпорт і збереження; до створення чи імпорту проєкту інші вкладки заблоковані. Вкладка «Конструктор» (рис.~\ref{fig:ui-constructor}) має ліворуч панель із розгортними секціями (кімнати, відображення, матеріали, двері та арки, електрика, світильники, предмети на стіні), праворуч~-- тривимірну сцену з панеллю режимів «Огляд», «Перемістити», «Проводка».

\setscreenshot[fig:ui-constructor]{0.95\textwidth}{assets/ui-constructor.png}{Вкладка «Конструктор»: контрольна сцена з двох кімнат}

Вкладка «Проводка» (рис.~\ref{fig:ui-wiring}) містить параметри нового проводу (вузли, тип кабелю, кількість жил, діаметр), прапорець «через стелю», кнопку контрольної точки та список проводів із діями «Паралельний» і «Видалити».

\setscreenshot[fig:ui-wiring]{0.95\textwidth}{assets/ui-wiring.png}{Вкладка «Проводка»: маршрути проводів між елементами}

Вкладка «Розрахунки» (рис.~\ref{fig:ui-calc}) виводить результати по кімнатах і блоки: гіпсокартон, рейки на стелю, плінтуси, детальна довжина проводки та розкладка ГКЛ. Розкладку можна показати безпосередньо на 3D-моделі (рис.~\ref{fig:ui-drywall}): плити зафарбовано за статусом, кожна має номер, плита під курсором підсвічується, а відрізки викладаються поруч із кімнатою.

\setscreenshot[fig:ui-calc]{0.95\textwidth}{assets/ui-calc.png}{Вкладка «Розрахунки»: результати та блок «Гіпсокартон»}

\setscreenshot[fig:ui-drywall]{0.95\textwidth}{assets/ui-drywall.png}{Розкладка листів ГКЛ на 3D-моделі та відрізки поруч із кімнатою}

Взаємодія у 3D ґрунтується на променях: вибір об'єкта виконується трасуванням променя від курсору до сцени з порогом 0,12~м для ліній, щоб тонкі проводи легко захоплювалися; під час перетягування об'єкта перетинається площина кожної стіни та стелі всіх кімнат, а обирається найближчий перетин, що дозволяє перемістити об'єкт на іншу стіну чи кімнату. Гарячі клавіші: \texttt{Ctrl+Z}~-- скасувати, \texttt{Ctrl+S}~-- зберегти.

\customSubSection[sec:verification]{Верифікація, продуктивність і обмеження}{1em}

\customSubSubSection[sec:verification-scene]{Контрольна сцена}{0.5em}

Для перевірки побудовано контрольну сцену: вітальня $5\times4\times2{,}7$~м і кухня $4\times4\times2{,}7$~м із підлогою на 0,15~м вище; двері на південній стіні вітальні (відкіс 0,2~м) та арка у спільній стіні з боку кожної кімнати (відкоси 0,1~м, у кухні~-- поріг 0,15~м); 11 елементів електрообладнання, 3 світильники, 2 предмети на стінах та 6 проводів трьох типів. Сцену завантажено в застосунок через імпорт, а показники зчитано з вкладки «Розрахунки».

\customSubSubSection[sec:verification-indep]{Незалежний перерахунок}{0.5em}

Незалежність забезпечено тим, що перерахунок виконує окремий скрипт мовою Python, який читає той самий файл проєкту, але не використовує жодного рядка коду застосунку. Результати порівняно в \tabref{tab:verify}: всі показники збіглися.

\begin{spacing}{1.15}
\small\noindent
\begin{xltabular}{\textwidth}{| >{\raggedright\arraybackslash}X | c | c | c |}
    \caption{Порівняння результатів застосунку з незалежним перерахунком} \label{tab:verify} \\
    \hline
    \textbf{Показник} & \textbf{Застосунок} & \textbf{Скрипт} & \textbf{Збіг} \\ \hline
    \endfirsthead

    \hline
    \textbf{Показник} & \textbf{Застосунок} & \textbf{Скрипт} & \textbf{Збіг} \\ \hline
    \endhead

    Чиста площа стін, вітальня / кухня, м$^2$ & 44,82 / 41,31 & 44,82 / 41,31 & так \\ \hline
    Площа прорізів разом, м$^2$ & 5,67 & 5,67 & так \\ \hline
    Площа відкосів, м$^2$ & 2,04 & 2,04 & так \\ \hline
    Листів ГКЛ за площею (стіни) & 33 & 33 & так \\ \hline
    Профіль UD, м / планок & 34,00 / 12 & 34,00 / 12 & так \\ \hline
    Профіль CD, м / планок & 63,00 / 21 & 63,00 / 21 & так \\ \hline
    Підвіси / краби, шт & 112 / 112 & 112 / 112 & так \\ \hline
    Підлоговий плінтус, м / планок & 31,30 / 14 & 31,30 / 14 & так \\ \hline
    Стельовий плінтус, м / планок & 34,00 / 15 & 34,00 / 15 & так \\ \hline
    Провід 1 (щиток~$\to$ коробка), м & 6,00 & 6,00 & так \\ \hline
    Провід 2 (коробка~$\to$ розетка), м & 3,10 & 3,10 & так \\ \hline
    Провід 3 (вимикач~$\to$ світильник), м & 4,69 & 4,69 & так \\ \hline
    Провід 4 (щиток~$\to$ кухня), м & 9,55 & 9,55 & так \\ \hline
    Провід 5 (щиток~$\to$ розетка, UTP), м & 7,10 & 7,10 & так \\ \hline
    Провід 6 (щиток~$\to$ вивід, оптика), м & 7,70 & 7,70 & так \\ \hline
    Сумарна довжина / із запасом 15~\%, м & 38,14 / 43,86 & 38,14 / 43,86 & так \\ \hline
\end{xltabular}
\end{spacing}

Ручна перевірка першого проводу за \eqref{eq:wirelen}: щиток розташований на висоті 1,5~м, коробка на висоті 2,4~м, рівень прокладання $c_y=2{,}7-0{,}02=2{,}68$~м, зсуви $|\Delta x|=0{,}6$~м і $|\Delta z|=3{,}94$~м; $L=1{,}18+0{,}6+3{,}94+0{,}28=6{,}00$~м, що збігається з виведеним значенням.

\customSubSubSection[sec:verification-dw]{Розкладка ГКЛ: порівняння режимів}{0.5em}

Нижня межа кількості листів \eqref{eq:lowerbound} для контрольної сцени (чиста площа з відкосами 88,17~м$^2$) становить $\lceil 88{,}17/3{,}0\rceil=30$, а проста оцінка за \eqref{eq:sheets}~-- 33. Результати розкладки за різних параметрів наведено в \tabref{tab:dwmodes}. Назви режимів «Горизонтально» й «Вертикально» подано так, як в інтерфейсі; розмір плитки в них становить $1{,}2\times2{,}5$ та $2{,}5\times1{,}2$~м відповідно.

\begin{spacing}{1.15}
\small\noindent
\begin{xltabular}{\textwidth}{| >{\raggedright\arraybackslash}X | c | c | c | c | c | c |}
    \caption{Результати розкладки ГКЛ для стін контрольної сцени} \label{tab:dwmodes} \\
    \hline
    \textbf{Режим} & \textbf{Шахи} & \textbf{Зал.} & \textbf{Нових} & \textbf{З залишку} & \textbf{Зі шматків} & \textbf{Залишок, м$^2$} \\ \hline
    \endfirsthead

    \hline
    \textbf{Режим} & \textbf{Шахи} & \textbf{Зал.} & \textbf{Нових} & \textbf{З залишку} & \textbf{Зі шматків} & \textbf{Залишок, м$^2$} \\ \hline
    \endhead

    Авто & так & так & 31 & 31 & 3 & 0 \\ \hline
    Авто & ні & так & 31 & 20 & 6 & 0,81 \\ \hline
    Авто & так & ні & 34 & 31 & --- & 13,42 \\ \hline
    Авто & ні & ні & 37 & 20 & --- & 22,43 \\ \hline
    Горизонтально & так & так & 30 & 47 & --- & 0,98 \\ \hline
    Горизонтально & ні & так & 31 & 46 & --- & 3,90 \\ \hline
    Вертикально & так & так & 31 & 38 & 2 & 2,86 \\ \hline
    Вертикально & ні & так & 31 & 27 & 5 & 0,81 \\ \hline
    Вертикально & так & ні & 33 & 38 & --- & 10,06 \\ \hline
    Вертикально & ні & ні & 36 & 27 & --- & 18,99 \\ \hline
    Економія & --- & так & 30 & 45 & 3 & 0,70 \\ \hline
\end{xltabular}
\end{spacing}

Примітка до таблиці: «Зал.»~-- прапорець використання залишків; для режиму «Горизонтально» значення за вимкненого прапорця залишків збігаються з наведеними й тому не повторюються, так само для режиму «Економія»; шаховий порядок у режимі «Економія» не застосовується.

Спостереження. По-перше, використання залишків скорочує кількість нових листів із 34 до 31 у режимі «Авто» та з 33 до 31 у режимі «Вертикально» за шахового порядку. По-друге, режими «Горизонтально» та «Економія» купують 30 листів, що збігається з нижньою межею; отже, на цій сцені знайдений розв'язок оптимальний за кількістю листів. По-третє, відмова від шахового порядку в режимі «Авто» зі залишками дає ті самі 31 лист, а без залишків~-- 37. По-четверте, найкращий режим заощаджує 3 листи (9~\%) порівняно з оцінкою за площею; з урахуванням стелі (площа 36~м$^2$) оцінка дає $33+14=47$ листів, а розкладка в режимі «Авто» зі залишками~-- 43. Ці результати отримано для однієї сцени та одного розміру листа й не узагальнюються на інші проєкти.

Внутрішню узгодженість перевірено інваріантом: кількість нових листів дорівнює кількості плит зі статусами «ціла», «ціла з отвором» та «розріз», а загальна кількість плит дорівнює сумі цих плит, плит із залишку та плит зі шматків. Інваріант виконано для всіх рядків \tabref{tab:dwmodes} (наприклад, для режиму «Авто» зі шматками: $12+6+13=31$ нових листів, $31+31+3=65$ плит). Окремий запуск модуля \texttt{drywall.js} у браузері на еквівалентних поверхнях дав ті самі кількості нових листів (31, 30, 31 і 30 для режимів «Авто», «Горизонтально», «Вертикально» й «Економія»; 43 із урахуванням стелі).

\customSubSubSection[sec:verification-perf]{Продуктивність}{0.5em}

Час розрахунку розкладки виміряно як середнє за 50 запусків модуля \texttt{drywall.js} у браузері на контрольній сцені (\tabref{tab:perf}). Усі значення нижчі за поріг 0,5~с із вимоги NFR-01; найповільнішим є режим «Горизонтально» (імовірно, через повторні спроби складання зі шматків у проході~3).

\begin{spacing}{1.15}
\small\noindent
\begin{xltabular}{\textwidth}{| >{\raggedright\arraybackslash}X | c | c |}
    \caption{Час розрахунку розкладки ГКЛ (середнє за 50 запусків)} \label{tab:perf} \\
    \hline
    \textbf{Режим} & \textbf{Час, мс} & \textbf{Нових листів} \\ \hline
    \endfirsthead

    \hline
    \textbf{Режим} & \textbf{Час, мс} & \textbf{Нових листів} \\ \hline
    \endhead

    Авто & 0,05 & 31 \\ \hline
    Горизонтально & 68,32 & 30 \\ \hline
    Вертикально & 0,92 & 31 \\ \hline
    Економія & 0,22 & 30 \\ \hline
    Авто, зі стелею & 1,62 & 43 \\ \hline
\end{xltabular}
\end{spacing}

\customSubSubSection[sec:limits]{Обмеження поточної версії}{0.5em}

За результатами аналізу коду та експериментів виявлено такі обмеження.

\begin{itemize}[noitemsep, topsep=0.3em, leftmargin=*]
    \item \textbf{Відсутні електротехнічні розрахунки.} Застосунок моделює розміщення й трасування та визначає довжини проводів, але не обчислює навантаження, переріз жил, падіння напруги, струми короткого замикання та параметри захисту. Поля «кількість жил» і «діаметр» зберігаються й виводяться в списку, але в розрахунках не використовуються. Це основний напрям подальшої роботи.
    \item \textbf{Немає автоматизованих тестів.} Коректність підтверджена одноразовим незалежним перерахунком на одній сцені; доцільно оформити його як набір регресійних тестів.
    \item \textbf{Евристичний розкрій.} Алгоритм не гарантує оптимуму; прорізи прямокутні (заокруглення арки ігнорується); результат залежить від прапорців і розміру листа.
    \item \textbf{Позначення режимів розкладки.} Назви «Горизонтально» та «Вертикально» в інтерфейсі відповідають розміру плитки (ширина~$\times$~висота), а не напрямку довгої сторони листа; для користувача це може бути неочевидним і потребує уточнення підписів.
    \item \textbf{Безпека введення.} Назви кімнат підставляються в HTML без екранування; для локального застосунку це прийнятно, але імпорт недовіреного файла потенційно небезпечний. Рекомендується екранувати текст або використовувати \texttt{textContent}.
    \item \textbf{Залежність від мережі та браузера.} Three.js завантажується з CDN; перезапис відкритого файлу на місці працює лише в Chromium~\citecustom{24}.
    \item \textbf{Спрощення розрахунків матеріалів.} Профілі каркаса обчислюються за фіксованими припущеннями (крок підвісів 0,6~м, кожна кімната окремо), плінтус віднімає ширину всіх прорізів, довжина проводів не враховує вертикальні заглиблення в стіни понад відсоток запасу.
\end{itemize}

\customSubSection[sec:ch2-conclusions]{Висновки до розділу II}{1em}

\begin{enumerate}[noitemsep, topsep=0.3em, leftmargin=*, label=\arabic*.]
    \item Обґрунтовано вибір клієнтського застосунку без збірки: профіль задачі не потребує сервера, а вимоги приватності та простоти розгортання задовольняються статичними файлами, ES-модулями та Three.js.
    \item Описано архітектуру за C4, модель даних проєкту у форматі JSON версії~2 із нормалізацією під час імпорту та історією скасування на 60 кроків.
    \item Формалізовано геометричну модель приміщення зі стіновою системою координат та алгоритм ортогонального трасування проводів із довжиною за формулою \eqref{eq:wirelen}.
    \item Виведено формули розрахунків площ, листів, профілів, плінтусів і проводки та описано алгоритм розкрою ГКЛ із пулом відрізків, гільйотинним розрізом, складанням зі шматків і режимом «Економія».
    \item Результати застосунку на контрольній сцені збіглися з незалежним перерахунком за всіма 16 показниками; режими «Горизонтально» й «Економія» досягли нижньої межі 30 листів проти 33 за простою оцінкою; час розкладки не перевищив 68~мс.
    \item Визначено обмеження версії, головним з яких є відсутність електротехнічних розрахунків і автоматизованих тестів.
\end{enumerate}
